\documentclass[../main]{subfiles}
\begin{document}
\chapter{Marco Jurídico de Proyectos}
\img{images/08/constitucion.jpg}{Constitución Argentina.}
\section{Pirámide Jurídica}
\img{images/08/piramide.jpg}{Pirámide Jurídica}
\info{La supremacía constitucional}{
  La supremacía constitucional es un principio teórico del Derecho
  constitucional que postula, originalmente, ubicar la Constitución
  de un país jerárquicamente por encima de todo el ordenamiento
  jurídico de ese país, considerándola como Ley Suprema del Estado y
  fundamento del sistema jurídico. Según cada país los tratados
  internacionales, convenciones internacionales o pactos
  internacionales gozarán o no del mismo rango que la Constitución
  Nacional de cada uno de ellos.

  El principio de supremacía constitucional es parte de un principio
  más general del derecho, llamado principio de jerarquía. Este
  principio tiene dos funciones: una positiva, de fundamentar lo
  inferior, y otra negativa, de hacer caer, anular o dejar sin
  efectos a aquello inferior que se le oponga.
}

\section{Legalidad de un proyecto}
Ningún proyecto infringir la ley, inmediatamente que un proyecto
incumple la ley, se transforma en inviable técnicamente, y debe ser
descartado o adaptado para que no infrinja ninguna ley o norma nacional.
\subsection{Principio de irretroactividad de la ley}
La irretroactividad es el fenómeno que produce que las normas no
tengan efectos hacia atrás en el tiempo. De esta manera se asegura
que dichos efectos comiencen en el momento de su entrada en vigor,
con la finalidad de dotar al ordenamiento jurídico de seguridad. Por
tanto un proyecto puede ser legal o viable hoy, pero mañana puede
dejar de ser, por una ley nueva.

\section{Instituto Nacional de la Propiedad Industrial}

Es un ente autárquico en el ámbito del ministerio de desarrollo
productivo y es la autoridad de aplicación de las leyes de protección
de los derechos de propiedad industrial.

\actividad{ Investigar y contestar:}
\begin{enumerate}
  \item ¿Qué es la propiedad intelectual?
  \item ¿Dónde puede registrarse la propiedad intelectual en Argentina?
  \item ¿Cuáles son los alcances de la protección?
  \item ¿Qué es una marca?
  \item ¿Qué es un dominio de internet?
\end{enumerate}

\end{document}